\documentclass{article}
\usepackage[T1]{fontenc}
\usepackage{lmodern}
\usepackage{graphicx}
\usepackage{amsmath,amssymb}
\usepackage[a4paper,margin=2cm]{geometry}
\usepackage[absolute,overlay]{textpos}
\setlength{\TPHorizModule}{1mm}
\setlength{\TPVertModule}{1mm}
\usepackage[indonesian]{babel}
\usepackage{hyperref}
\setlength{\emergencystretch}{2em}
% unit: Halaman 20

\begin{document}

% === Halaman 20 (llm) ===

\textbf{Contoh 2:} Misalkan $a = 01010111 = x^6 + x^4 + x^2 + x + 1$ dan $b = 10000011 = x^7 + x + 1$\
$a$ dan $b \in \text{GF}(2^8)$ (dalam notasi Hex, $a = \text{'57'}$ dan $b = \text{'83'}$)\
(i) $a + b = \text{'57'} + \text{'83'} = (x^6 + x^4 + x^2 + x + 1) + (x^7 + x + 1)$\
\begin{align*}$\\ = x^7 + x^6 + x^4 + x^2 + x + 2 \pmod{2} \quad \text{\color{red}Bagi tiap koefisien dengan 2,}\\ &\text{\color{red}lalu ambil sisanya}\\ = x^7 + x^6 + x^4 + x^2 + x + 0$\\ $= x^7 + x^6 + x^4 + x^2 + x = 11010100 = \text{'D4'}$\end{align*}\
Dalam representasi biner:
\begin{center}
\begin{tabular}{rl}
01010111 \\ $\underline{10000011} +$ \\ 11010100 $\rightarrow$ sama dengan hasil operasi XOR
\end{tabular}
\end{center}
$\therefore a + b = 11010100 = a \text{ XOR } b$

\vspace{10mm}
\begin{center}
\small Rinaldi Munir/II4021 Kriptografi
\end{center}

\end{document}
